\documentclass[11pt]{article}
\usepackage[margin=1in]{geometry}
\usepackage{amsmath,amssymb,amsthm}
\usepackage{hyperref}
\usepackage{enumitem}
\usepackage{xurl}

\newtheorem{theorem}{Theorem}
\newtheorem{lemma}[theorem]{Lemma}
\newtheorem{corollary}[theorem]{Corollary}
\theoremstyle{definition}
\newtheorem{definition}[theorem]{Definition}
\theoremstyle{remark}
\newtheorem{remark}[theorem]{Remark}

\let\oldus\_
\renewcommand{\_}{\oldus\allowbreak}
\newcommand{\C}{\mathbb{C}}

\title{A disproof of Conjecture 00000002167\\[2pt]
\large The spectral-radius maximizer among graphs of girth $\ge 5$ is not unique}
\date{}

\begin{document}
\sloppy
\maketitle

\section{The conjecture}

Conjecture 00000002167 reads:

\begin{quote}
\emph{Definition:} The spectral radius of $n$-vertex graphs of girth $\ge 5$. \emph{Conjecture:} The spectral
radius is attained by Moore-approximate graphs (Brown type or PC-lab type); the classification of extremal
graphs holds for all $n\ge5$ (uniqueness of spectral-radius extrema).
\end{quote}

\noindent The statement is a conjunction. Its last clause, made explicit by the parenthetical, is:
\begin{quote}
(U) for every $n\ge5$, the graphs that attain the maximum spectral radius among $n$-vertex graphs of girth
$\ge5$ are unique up to isomorphism.
\end{quote}
The terms ``Moore-approximate'', ``Brown type'' and ``PC-lab type'' are not defined in the statement, and we
do not need them: we show that (U) is false at $n=5$, so the conjunction is false whatever the first clause
means. The failure persists if only connected graphs are considered.

\section{Definitions}

\begin{definition}
The \emph{girth} of a graph is the length of a shortest cycle; a graph with no cycle (a forest) has girth
$\infty$. ``If the graph does not contain any cycles (that is, it is a forest), its girth is defined to be
infinity'' (Wikipedia, \emph{Girth (graph theory)},
\url{https://en.wikipedia.org/wiki/Girth_(graph_theory)}). In Lean this is Mathlib's
\texttt{SimpleGraph.egirth} (value in $\mathbb N\cup\{\infty\}$, equal to $\top$ exactly for acyclic graphs),
and ``girth $\ge5$'' is \texttt{5 <= G.egirth}.
\end{definition}

\begin{definition}
The \emph{spectral radius} $\rho(G)$ of a finite graph is $\max\{|\lambda| : \lambda$ an eigenvalue of its
adjacency matrix$\}$. In Lean, \texttt{adjSpectralRadius G} is Mathlib's
\texttt{spectralRadius} $\C$ \texttt{(G.adjMatrix} $\C$\texttt{)}, i.e.\ the supremum of $\|\lambda\|$ over
the spectrum of the complex adjacency matrix.
\end{definition}

\noindent $C_5$ is Mathlib's \texttt{cycleGraph 5} (vertex $i$ adjacent to $i\pm1 \bmod 5$). $K_{1,4}$ is
\texttt{starGraph 4} on $\{0,\dots,4\}$: $i\sim j$ iff $i\ne j$ and $0\in\{i,j\}$. $C_5$ is the smallest
girth-$5$ (diameter-$2$) Moore graph: Wikipedia (\emph{Moore graph}) lists ``C5 with diameter 2, girth 5, degree 2, and order 5''.

\section{The bound $\rho(G)\le\sqrt{n-1}$ for girth $\ge5$}

\begin{lemma}\label{lem:girth}
Let $G$ have girth $\ge5$. Then $G$ has no triangle $a\sim b\sim c\sim a$, and no $4$-cycle
$a\sim b\sim c\sim d\sim a$ with $a\ne c$, $b\ne d$. Conversely, a graph with neither has girth $\ge5$.
\end{lemma}

\begin{proof}
A triangle or such a quadrilateral is a cycle of length $3$ or $4$. Conversely, every cycle has length
$\ge3$, and a cycle of length $3$ or $4$ is a triangle or such a quadrilateral (its support is duplicate-free).
\end{proof}

\begin{lemma}\label{lem:deg}
If $G$ has girth $\ge5$ and $n$ vertices, then $\sum_{l\sim i}\deg l\le n-1$ for every vertex $i$.
\end{lemma}

\begin{proof}
For $l\sim i$ let $T_l=N(l)\setminus\{i\}$, so $|T_l|=\deg l-1$. If $j\in T_l\cap T_{l'}$ with $l\ne l'$, then
$i\sim l\sim j\sim l'\sim i$ is a $4$-cycle with $i\ne j$, $l\ne l'$. If $j\in T_l$ and $j\sim i$, then
$i,l,j$ is a triangle. So the $T_l$ are pairwise disjoint subsets of $V\setminus(\{i\}\cup N(i))$, and
$\sum_{l\sim i}(\deg l-1)\le n-1-\deg i$, i.e.\ $\sum_{l\sim i}\deg l\le n-1$.
\end{proof}

\begin{theorem}\label{thm:bound}
If $G$ has girth $\ge5$ and $n$ vertices, every eigenvalue $\lambda$ of its adjacency matrix $A$ satisfies
$|\lambda|^2\le n-1$; hence $\rho(G)\le\sqrt{n-1}$.
\end{theorem}

\begin{proof}
An element of the spectrum has an eigenvector $v\ne0$ with $Av=\lambda v$ (the determinant of
$\lambda I-A$ vanishes). Choose $i$ with $|v_i|$ maximal; then $|v_i|>0$. Since
$(Av)_l=\sum_{j\sim l}v_j=\lambda v_l$,
\[
  |\lambda|^2|v_i| = \Big|\sum_{l\sim i}\sum_{j\sim l}v_j\Big| \le \sum_{l\sim i}\deg(l)\,|v_i|
  \le (n-1)|v_i|
\]
by Lemma~\ref{lem:deg}. Divide by $|v_i|$.
\end{proof}

\section{Two non-isomorphic extremal graphs on five vertices}

\begin{theorem}[main]\label{thm:main}
\begin{enumerate}[label=(\alph*)]
  \item $C_5$ and $K_{1,4}$ have girth $\ge5$, and both are connected.
  \item $\rho(C_5)=\rho(K_{1,4})=2$, and $\rho(G)\le2$ for every graph $G$ of girth $\ge5$ on $5$ vertices.
  \item Hence both $C_5$ and $K_{1,4}$ attain the maximum spectral radius among $5$-vertex graphs of girth
        $\ge5$, and also among connected ones.
  \item $C_5\not\cong K_{1,4}$.
\end{enumerate}
\end{theorem}

\begin{proof}
(a) Neither graph contains a triangle or a quadrilateral as in Lemma~\ref{lem:girth} (a finite check over
$\{0,\dots,4\}$). $C_5$ is connected (Mathlib \texttt{cycleGraph\_connected}); every vertex of $K_{1,4}$ is
$0$ or adjacent to $0$. (b) The all-ones vector is an eigenvector of $C_5$ with eigenvalue $2$, since $C_5$
is $2$-regular. For $K_{1,4}$, $v=(2,1,1,1,1)$ gives $Av=(4,2,2,2,2)=2v$. So $2$ lies in both spectra and
$\rho\ge2$; Theorem~\ref{thm:bound} with $n=5$ gives $\rho\le\sqrt4=2$ for every girth-$\ge5$ graph on
$5$ vertices. (c) follows from (a) and (b). (d) An isomorphism preserves the number of edges, and $C_5$ has
$5$ edges while $K_{1,4}$ has $4$.
\end{proof}

\begin{corollary}
Clause (U) fails at $n=5$, both among all graphs and among connected graphs. Hence Conjecture 00000002167 is
false.
\end{corollary}

\begin{remark}
The same happens at every order of a girth-$5$ Moore graph: $n=10$ (Petersen graph and $K_{1,9}$) and
$n=50$ (Hoffman--Singleton graph and $K_{1,49}$), since in each case both graphs have spectral radius
$\sqrt{n-1}$. Only $n=5$ is formalized; one counterexample suffices for the ``for all $n\ge5$'' claim.
\end{remark}

\begin{remark}[reading]
The refutation uses the standard convention that forests have girth $\infty$, so that $K_{1,4}$ has girth
$\ge5$. With that convention the star $K_{1,n-1}$ attains the bound $\sqrt{n-1}$ for every $n$, so it is an
extremal graph for every $n$.
\end{remark}

\section{Lean 4 formalization}

The directory \texttt{lean/} is a Lean~4 project (Lean \texttt{v4.33.1}, Mathlib \texttt{v4.33.1}); the
source is \texttt{Conjecture2167/Basic.lean} (namespace \texttt{C2167}). It contains:
\begin{itemize}
  \item \texttt{adjSpectralRadius}, \texttt{starGraph}, and the predicates
        \texttt{IsSpecExtremal P G} (girth $\ge5$, side condition $P$, and maximal spectral radius among all
        graphs $H$ on \texttt{Fin n} with girth $\ge5$ and $P\,H$) and \texttt{UniqueExtremal n P}
        (any two such maximizers are isomorphic, \texttt{Nonempty (G $\simeq$g H)});
  \item \texttt{no\_triangle}, \texttt{no\_quad}, \texttt{five\_le\_egirth} (Lemma~\ref{lem:girth}),
        \texttt{sum\_degree\_neighbors\_le} (Lemma~\ref{lem:deg}),
        \texttt{norm\_sq\_le\_of\_mem\_spectrum} and \texttt{spectralRadius\_le\_sqrt}
        (Theorem~\ref{thm:bound}, for any finite vertex type);
  \item \texttt{main} (Theorem~\ref{thm:main}): $\rho(C_5)=2$, $\rho(K_{1,4})=2$, the bound on
        \texttt{Fin 5}, \texttt{IsSpecExtremal} for both graphs with $P=$ \texttt{True} and with
        $P=$ \texttt{Connected}, and \texttt{IsEmpty (cycleGraph 5 $\simeq$g starGraph 4)};
  \item \texttt{conjecture\_false}:
        $\neg(\forall n\ge5,\ \texttt{UniqueExtremal}\ n\ \texttt{True})\ \wedge\
        \neg(\forall n\ge5,\ \texttt{UniqueExtremal}\ n\ \texttt{Connected})$.
\end{itemize}
There is no \texttt{sorry}, no \texttt{native\_decide} and no added axiom. \texttt{\#print axioms} for
\texttt{C2167.conjecture\_false}, \texttt{C2167.main} and \texttt{C2167.spectralRadius\_le\_sqrt} reports
only \texttt{propext}, \texttt{Classical.choice} and \texttt{Quot.sound}. The file also compiles with
\texttt{-DwarningAsError=true}. To check:
\begin{verbatim}
cd lean
lake exe cache get
lake build
\end{verbatim}

\end{document}
